% Tabela-síntese da seção 2.2.2 (Características do mercado de criptoativos)
\SinteseInicio{Evidências empíricas sobre o mercado de criptoativos}{tab:sintese_cripto}

\SinteseAutor{charfeddine2020investigating} &
Criptoativos e ativos convencionais. &
Retornos diários de \textit{Bitcoin} e \textit{Ethereum} frente a \acs{sp500}, ouro e petróleo (2015--2018); correlações, cópulas e modelos GARCH; pesos ótimos de carteira. &
Desvios-padrão diários de 1,723\% (BTC) e 2,983\% (ETH), contra 0,239\% (\acs{sp500}), 0,265\% (ouro) e 0,744\% (petróleo); correlações de Pearson do BTC próximas de zero (0,010 com o \acs{sp500}; 0,017 com o ouro; \textminus0,008 com o petróleo); pesos médios ótimos de BTC de 4,9\% com o \acs{sp500} e 5,7\% com o ouro. &
Caracteriza a volatilidade e a baixa dependência dos criptoativos e sustenta seu potencial de diversificação em pequenas proporções. \\
\SinteseLinha

\SinteseAutor{wang2020stablecoins} &
Papel das \textit{stablecoins}. &
\textit{Stablecoins} atreladas ao dólar e ao ouro frente a criptomoedas tradicionais (2015--2019, dados diários); redução de \acs{var} e \acs{cvar} em carteiras. &
Em condições normais, as \textit{stablecoins} atuam como diversificadoras e só em circunstâncias específicas como porto seguro; o \textit{Tether} reduziu \acs{var} e \acs{cvar} em várias combinações, mas outras \textit{stablecoins} elevaram o risco em algumas configurações. &
Mostra que o efeito de proteção depende do ativo, da condição de mercado e da composição da carteira. \\
\SinteseLinha

\SinteseAutor{baur2021crypto} &
Proteção contra quedas do \textit{Bitcoin}. &
Variações de preço em alta frequência das maiores \textit{stablecoins} frente ao \textit{Bitcoin}. &
As \textit{stablecoins} funcionam como porto seguro em relação ao \textit{Bitcoin}, embora não sejam estáveis o tempo todo. &
Sustenta a volatilidade do \textit{Bitcoin} também em bases diárias e complementa a discussão sobre \textit{stablecoins}. \\
\SinteseLinha

\SinteseAutor{goodell2021diversifying} &
Diversificação durante crises. &
Correlações de BTC, ETH, LTC e \textit{Tether} com índices acionários (2019--2021), com subperíodo da pandemia de \acs{covid19}. &
Na pandemia, correlações superiores a 0,80 entre algumas das principais criptomoedas; o \textit{Tether} teve correlações negativas com o \acs{sp500} (\textminus0,230), \acs{ibex35} (\textminus0,186), \acs{dax} (\textminus0,156) e \acs{cac40} (\textminus0,153). &
Mostra que os benefícios de diversificação dependem do tipo de criptoativo e das condições de mercado. \\
\SinteseLinha

\SinteseAutor{fang2022cryptocurrency} &
Pesquisa sobre negociação de criptoativos. &
Revisão de 146 trabalhos (2013 a junho de 2021) sobre sistemas de negociação, bolhas e condições extremas, previsão de volatilidade e retorno, carteiras e análise técnica. &
Organiza a literatura da área e aponta decisões baseadas em informação limitada, horizonte de curto prazo e limitações de escalabilidade do mercado. &
Contextualiza os riscos do mercado e o uso crescente de \acs{ml} na negociação de criptoativos. \\

\SinteseFim
